\section{Preliminaries}
\label{sec:preliminaries}

\subsection{Markov decision processes and off-policy value learning}

\begin{equation}
Q^\star(s,a)= r(s,a) + \gamma\,\mathbb{E}_{s'\sim P(\cdot|s,a)}
\left[\max_{a'\in\mathcal{A}} Q^\star(s',a')\right].
\label{eq:bellman}
\end{equation}
Keep the RL background minimal here and focus to have space in the next section to talk about federation

\subsection{Hyperdimensional computing (HDC)}

\label{subsec:hdc_background}
{\color{red} DESCRIBE ONLY GENERAL HDC HERE COMPACT. }
Hyperdimensional computing (HDC) is a brain-inspired computational paradigm in which symbols and structured entities are represented as high-dimensional vectors---called \emph{hypervectors}---with components drawn independently from simple distributions~\cite{kanerva2009hyperdimensional}. In such spaces, independently sampled hypervectors are nearly orthogonal with high probability, a geometric property that underlies classical HDC operations such as \emph{bundling} (superposition via addition), \emph{binding} (association via elementwise multiplication or permutation), and \emph{permutation} (role shifting). Because information is distributed holographically across all dimensions, HDC representations exhibit strong robustness to noise, quantization, and partial corruption.

Given an input $x$, an HDC encoder produces a bounded hypervector $h=\phi(x)\in\mathbb{R}^D$ (or $\mathbb{C}^D$) using random projections or compositional schemes built from a small set of base hypervectors. These encoders admit a precise random-feature interpretation: the empirical kernel
\begin{equation}
\label{eq:empirical_kernel}
k_D(x,x')=\langle \phi(x),\phi(x')\rangle
\end{equation}
converges to a smooth limiting kernel $k_\ast(x,x')$ as $D\to\infty$, with uniform concentration at rate $O(D^{-1/2})$ \cite{rahimi2007random,bach2015equivalence,rudi2017generalization}. As a consequence, linear prediction in the hypervector domain serves as a computationally efficient, finite-dimensional approximation to kernel methods in the RKHS associated with $k_\ast$.



% \begin{enumerate}
%     \item \textbf{Reward heterogeneity}
%     \item \textbf{}
% \end{enumerate}


% \subsection{Environmental and Model-wise Heterogeneity}
% \label{subsec:hetero}
% We consider $N$ agents. Agent $i$ interacts with its own environment modeled by an MDP
% \begin{equation}
% \mathcal{M}^{(i)}=(\mathcal{S},\mathcal{A},P^{(i)},r^{(i)},\gamma),
% \end{equation}
% where all agents share the same state--action space $(\mathcal{S},\mathcal{A})$ but may differ in reward and transition dynamics.
% For notational convenience, we write $x=(s,a)\in\mathcal{S}\times\mathcal{A}$.

% We quantify environmental mismatch across agents using
% \begin{equation}
% \label{def:hetero}
% \begin{aligned}
% \varepsilon_r
% &:=
% \max_{i,j}\ \frac{\sup_{(s,a)}\big|r^{(i)}(s,a)-r^{(j)}(s,a)\big|}{R},
% \,\\
% \varepsilon_p
% &:=
% \max_{i,j}\ \sup_{(s,a)}
% d_{\mathrm{TV}}\!\big(P^{(i)}(\cdot\mid s,a),P^{(j)}(\cdot\mid s,a)\big), \\
% \varepsilon_m
% &:=
% \max_{i,j}\ 

% \end{aligned}
% \end{equation}
% where $d_{\mathrm{TV}}$ denotes total variation distance.\footnote{Other equivalent discrepancy measures can be used; we adopt total variation for concreteness.}, R denotes the 
% In the finite-time analysis, heterogeneity enters through an induced term $\Lambda(\varepsilon_p,\varepsilon_r)$, which captures the irreducible bias floor arising from aggregating update directions generated by non-identical environments.






